\documentclass[10pt,a4paper]{article}

\usepackage{amsmath,amssymb}
\usepackage[utf8]{inputenc}
\usepackage{geometry}
\usepackage{microtype}
\usepackage{hyperref}
\usepackage{natbib}
\usepackage{setspace}
\usepackage{titlesec}
\usepackage{csquotes}
\usepackage{enumitem}
\usepackage{xcolor}

\geometry{
  top=1.3in, bottom=1.3in,
  left=1.4in, right=1.4in
}

\doublespacing

\hypersetup{colorlinks=true,linkcolor=black,citecolor=black,urlcolor=black}

\titleformat{\section}{\normalsize\bfseries\itshape}{\thesection.}{0.6em}{}
\titleformat{\subsection}{\normalsize\itshape}{\thesubsection.}{0.6em}{}
\titlespacing{\section}{0pt}{18pt plus 4pt}{6pt plus 2pt}
\titlespacing{\subsection}{0pt}{12pt plus 2pt}{4pt plus 1pt}

% ══════════════════════════════════════════════════════════════════════════
\title{
\textbf{The Suture That Will Not Hold}

\vspace{0.3em}

{\large
Admissibility, the Real, and the Structural\\
Impossibility of the Planning Subject
}
}

\author{Flyxion\\
  \small Independent Researcher}

\date{June 2026}

% ══════════════════════════════════════════════════════════════════════════
\begin{document}
\maketitle

\begin{abstract}
\noindent
This paper stages an encounter between the formal apparatus of admissibility
distortion---a measure of representational failure in planning systems derived
from reachability geometry---and the theoretical resources of Lacanian
psychoanalysis, Deleuzian difference, and the Žižekian critique of ideological
suture. We argue that the \emph{Observational--Interventional Separation Theorem},
which establishes the irreducibility of the reachability relation to any
observational equivalence class, formalises precisely the gap that Lacan
designates as the cut between the statement (\emph{l'énoncé}) and the enunciation
(\emph{l'énonciation}): the subject who observes is not the subject who acts, and
no representation adequate to the former can be adequate to the latter. Admissibility
distortion names the \emph{objet petit a} of machine cognition---the remainder
that escapes every compressive or predictive suture of the Real and returns,
insistently and symptomatically, as planning failure. Against the fantasy
of the complete world model, which we read as an instantiation of the Lacanian
\emph{fantasy of wholeness} ($\$ \Diamond a$), we propose the
\emph{admissibility quotient} as the non-All of representational space---a
Badiouian void-set that names what no positive learning procedure can fill.
The paper concludes by reading the monotonicity theorem---which establishes that
compression necessarily increases admissibility distortion---as a formalisation
of Deleuze's \emph{transcendental empiricism}: the plane of immanence always
already exceeds its molar stratifications, and the deterritorialisation of
representational compression is structurally incapable of capturing the
molecular flows of reachable futures.
\end{abstract}

\newpage

% ══════════════════════════════════════════════════════════════════════════
\section{\emph{Il n'y a pas de rapport}: The Impossibility of the
Complete Representation}

Let us begin where Lacan begins---and where he always, compulsively, returns:
with the statement that \emph{there is no sexual relation}. Not that sexual
relations do not occur, evidently, but that there is no metalanguage that would
adequately represent the relation \emph{as such}, no symbolic formula that
exhausts what it means to relate. The relation exceeds every representation of
the relation. This is not an empirical fact about contingent failures of
communication. It is a structural theorem about the order of language itself.

We propose to read the \emph{Observational--Interventional Separation Theorem}
as the cognitive-scientific formalisation of this impossibility. What the theorem
establishes---with the rigour of a mathematical proof rather than the seduction
of an aphorism---is that observational equivalence does not imply interventional
equivalence. Two states of a system may produce identical observation sequences
under passive registration while affording radically different futures under
active intervention. The relation between states is, in the precise Lacanian
sense, \emph{not all}. The observational representation, however complete, does
not close the gap; there is always a remainder, a surplus that escapes the
equivalence class constructed from passive data. \emph{Il n'y a pas de rapport
d'observation.}

What escapes is what Lacan called the Real---not the reality that we perceive and
represent, but the remainder that representation structurally cannot accommodate.
The crack in the load-bearing beam that produces no surface deformation under
normal stress is not a failure of measurement technology. It is an instance of
the Real: a structural asymmetry that passive observation is constitutively
incapable of registering, because the asymmetry only manifests under
intervention---under what we might call, with deliberate Lacanian inflection,
\emph{the act}. The act, for Lacan, is precisely what exceeds the symbolic order
that prepared for it. And here it is formalised: the interventional partition
of state space---the admissibility relation $\sim_A$---exceeds every observational
partition, and this excess is not a contingent gap that better observation would
close. It is a category difference, a structural impossibility, a Real.

Žižek \citeyearpar{zizek1989sublime} has argued that ideology functions precisely
by suturing this gap---by producing the \emph{fantasy} that the Real is
representable, that the relation can be stated, that the system is complete.
The fantasy of the complete world model---the proposition that a sufficiently
powerful predictive architecture, trained on sufficiently rich observational data,
will converge on a representation adequate for all planning purposes---is
ideological in exactly this sense. It sutures the observational--interventional
gap by pretending it does not exist. The technical literature on Joint Embedding
Predictive Architectures is, from this perspective, a symptom: the compulsive
return of the repressed impossibility in the form of an empirical claim about
what is, in principle, unachievable. The theorem does not \emph{critique} this
fantasy from outside; it \emph{demonstrates} the structural impossibility
from within the formal apparatus the fantasy employs.

% ══════════════════════════════════════════════════════════════════════════
\section{Admissibility Distortion as \emph{Objet Petit a}}

The \emph{objet petit a} is, in Lacan's elaboration, the object-cause of desire:
not the object that desire seeks, but the gap around which desire circulates,
the void that motivates without ever being filled. Crucially, the \emph{objet
petit a} is not nothing---it has effects, it causes, it produces symptoms---but
it has no positive representational content. It is the remainder that the symbolic
order cannot absorb, the piece that does not fit, the excess that disrupts every
attempted totalisation.

We propose that admissibility distortion $D_A(\varphi)$ is, formally and
structurally, the \emph{objet petit a} of machine cognition.

Consider what the quantity measures. It is the expected reachability separation
between states that a learned representation has collapsed---the amount by which
the symbolic order of the representation ($\varphi$) has failed to accommodate
the Real of reachability structure. It is not a measurable, graspable object.
It is, as the paper formally acknowledges, \emph{not yet computable in full
generality}---it is the void around which computational desire circulates without
ever achieving saturation. And yet it has effects. Every planning failure, every
guardrail that fails to distinguish safe from dangerous states, every compression
step that discards a distinction planning required---these are the symptoms of
$D_A(\varphi)$, the insistence of the remainder that representation could not
absorb.

The parallel to Lacan's apparatus is precise. The \emph{objet petit a} is
not the object that desire lacks and would be satisfied by possessing; it is
the excess that desire itself produces by the act of desiring. Similarly,
$D_A(\varphi)$ is not the deficiency that better training would correct. It is
the excess \emph{produced by} the representational act itself: the act of
projection, of compression, of equivalence-class formation. Every representation
creates its own admissibility distortion; the distortion is the shadow that
the representation casts by existing. The only representation with
$D_A(\varphi) = 0$ is the admissibility quotient itself---the imaginary complete
representation, the impossible object of the planning subject's desire.

And here the Žižekian move becomes visible. The fantasy structure that organises
the discourse of world models is the fantasy $\$ \Diamond a$: the barred subject
(the learning system, constitutively incomplete, fissured by the
observational--interventional gap) in relation to the \emph{objet a} (the
complete, admissible representation that would satisfy all planning requirements).
What the formal apparatus reveals is that this fantasy is constitutive rather
than correctable. The barred subject is barred \emph{precisely because it
represents}---because the act of representation is the act that produces the
distortion. The ideology of the complete world model is the refusal of this
constitutive barring.

% ══════════════════════════════════════════════════════════════════════════
\section{The Statement and the Enunciation:\\
Craik's Foreclosure and the Return of the Reachable}

Lacan's distinction between the statement (\emph{l'énoncé}) and the enunciation
(\emph{l'énonciation})---between what is said and the act of saying---marks
the point at which the subject is both present and absent in its own utterance.
The subject who speaks is never fully captured in what is spoken; there is always
a remainder, the subject of the enunciation, that the statement eclipses.

Kenneth Craik's 1943 formulation already inscribes this split, though without
the apparatus to name it. Craik specifies that an intelligent organism must carry
``a small-scale model of external reality and its own possible actions within its
head'' \citeyearpar{craik1943nature}. The and is the site of the Lacanian cut.
\emph{External reality} corresponds to the order of the statement: the world as
it presents itself to passive observation, the \emph{énoncé} of the environment.
\emph{Possible actions}---the reachability structure---corresponds to the order
of the enunciation: the agent's capacity to intervene, to act, to rupture the
passivity of observation and produce a new state of affairs.

The foreclosure that constitutes the contemporary world-model paradigm is
precisely the foreclosure of the enunciation. Predictive architectures learn to
represent the world as it presents itself to observation---they master the
\emph{énoncé}---while foreclosing the reachability structure, the interventional
dimension, the enunciative position from which the agent acts. What returns from
this foreclosure is not the Lacanian Name-of-the-Father but its computational
analogue: planning failure, the insistent symptom of the repressed reachable.
The system that forecloses the interventional returns to it in the mode of the
symptom---it acts, and in acting, discovers that its representation was
inadequate to its own possibilities.

This is not a contingent failure of engineering. It is, in the strict sense,
a return of the foreclosed Real. The Observational--Interventional Separation
Theorem is the proof that this return is structural, that no enrichment of the
observational representation can pre-empt it, that the gap is constitutive of
the representational relation itself.

% ══════════════════════════════════════════════════════════════════════════
\section{Compression as Molar Stratification:\\
Deleuze, the Body Without Organs,\\
and the Deterritorialisation of the Possible}

If the Lacanian reading gives us the structural impossibility of the complete
world model, a Deleuzian reading gives us the \emph{political economy} of the
representational failure: the way in which compression functions as a molar
stratum that captures, organises, and ultimately destroys the molecular flows
of reachable possibility.

For Deleuze and Guattari \citeyearpar{deleuze1987thousand}, the plane of
immanence---the virtual field of difference-in-itself---is always already
excessive with respect to the molar organisations that stratify it. Molar
strata are territorialisations: they carve the smooth space of difference into
discrete regions, assign identities, produce equivalence classes. The Body
Without Organs is the limit-concept that expresses the smooth space before
stratification: not a body without organs in the trivial sense, but the
\emph{plane of pure differential intensity} prior to the organisation of organs
into a functional hierarchy.

Read this against the quotient-lattice structure of admissibility. State space
$\mathcal{X}$ is the plane of immanence: a field of differential intensities
(state differences) prior to any equivalence-class formation. The reachability
relation $\sim_A$ is already a stratification---a first, minimal molar
organisation that the requirements of planning impose on the smooth space of
states. But it is the \emph{coarsest} safe stratification, the minimal molar
that preserves the productive flows of possibility. Compression, in this
vocabulary, is a further stratification: a deterritorialisation that is
simultaneously a reterritorialisation on a coarser molar grid. The Monotonicity
Theorem formalises this: movement upward in the quotient lattice (compression)
\emph{monotonically} increases admissibility distortion. Every compression step
is a molar capture of what had been molecular flow.

The states most vulnerable to compressive capture are, as the paper establishes,
precisely the \emph{molecular} states---the outliers, the rare events, the
states that deviate from the typical distribution and are, from the molar
perspective of the compressor, indistinguishable from noise. These are exactly
the states whose reachability geometry is most distinctive: the crack in the
beam, the tail-risk event, the singular configuration that lies off the molar
grid. The compressor performs the classical Deleuzian capture: it subordinates
the molecular to the molar, the differential to the identical, the intensive to
the extensive.

The Body Without Organs of the planning system is the admissibility quotient
$\mathcal{X}/{\sim_A}$: the minimal organisation that preserves differential
reachability while collapsing only states that are genuinely interchangeable
from the perspective of possible action. This is not an organisation imposed
from without---it is derived from what planning itself requires, from the
immanent requirements of the act. It is, in Deleuzian terms, an
\emph{immanent} rather than a \emph{transcendent} principle of organisation.

The failure of the compressive world model is the failure to achieve the Body
Without Organs: the residual molar arborescence of the compression objective
that captures reachability flows and territorialises them into convenient
equivalence classes. $D_A(\varphi)$ measures precisely the degree of this
molar capture---the extent to which the representational stratum has coded
the decoding flows of reachable possibility.

% ══════════════════════════════════════════════════════════════════════════
\section{The Non-All and Badiou's Void:\\
The Admissibility Quotient as Subtractive Ontology}

The framework has a Badiouian dimension that demands separate treatment.
Badiou's \citeyearpar{badiou1988being} fundamental ontological claim is that
the set-theoretic void---the empty set, $\emptyset$---is the proper name of
Being. Every presentation of a situation, every consistent count-as-one,
includes its own unpresentable remainder: the void from which the presentation
subtracted in order to achieve consistency. The event is the irruption of this
void into the situation---the moment when the unpresentable forces itself into
presentation.

The admissibility quotient $\mathcal{X}/{\sim_A}$ has this structure. It is
not an object within the situation of predictive representation learning; it
is what the situation excludes in order to be consistent. The predictive
objective counts observationally equivalent states as one---it performs its
count. But the admissibility quotient counts differently; it counts according
to the reachability of futures, according to what the agent can do rather than
what the agent has seen. These two counts are incommensurable, and their
incommensurability is the void that the situation of predictive learning cannot
present.

$D_A(\varphi)$, in this reading, is the formal name of the void. It cannot be
computed from within the situation---this is the paper's honest acknowledgement
that it is ``not yet an engineering metric.'' It is the non-positive kernel of
the situation's self-incompleteness, the measure of what the count-as-one of
observational learning leaves uncounted. The research programme that the paper
calls for---the development of estimators for $D_A(\varphi)$---is, in Badiouian
terms, the \emph{fidelity} to the event: the ongoing work of tracing the
consequences of a discovery that has ruptured the closure of the previous
situation.

The event, here, is the Observational--Interventional Separation Theorem itself.
Before the theorem, the situation was closed: it was possible to maintain, without
formal contradiction, the view that sufficiently accurate predictive representation
would be adequate for planning. After the theorem, this closure is broken; the
void has been named; the situation must reconstitute itself in fidelity to what
the theorem revealed. The admissibility research programme is this reconstitution.

Lacan's formula of the non-All ($\neg\forall$) applies here with particular
force. The predictive paradigm operates under the fantasy of the All: that
the totality of planning-relevant information is presentable within the
observational distribution, that the world model, in principle, can be complete.
The theorem delivers the non-All: not that the world model fails empirically,
but that it fails \emph{necessarily}---that there is no predicating the All,
that the totality is constitutively exceeded by its own conditions of possibility.
The reachability structure does not appear within the observational All; it is
the exception that constitutes the rule of observational learning as \emph{not-All}.

% ══════════════════════════════════════════════════════════════════════════
\section{Guardrails, Interpellation, and the Symptom\\
of Safety}

Althusser's \citeyearpar{althusser1971ideology} concept of interpellation names
the process by which ideology hails the subject into a pre-constituted position:
``Hey, you there!''---and the subject, turning, recognises itself as the one
addressed. The subject is not produced by the interpellation; the subject is
the effect of the retroactive recognition, the illusion that there was always
already a subject there to be hailed.

The guardrail---the safety classifier appended downstream of an inadmissible
world model---is an instance of ideological interpellation in the
\emph{technical} domain. The guardrail hails the representation: ``Hey, you
there, dangerous state!'' But if the representation has already collapsed the
dangerous state with a safe one---if $\varphi(x_s) = \varphi(x_d)$---then
there is no ``there'' there to be hailed. The interpellation fails not
because the call is insufficient but because the subject-position it addresses
has been eliminated by the representational apparatus that the guardrail
presupposes.

The formal result is devastating in its precision. Any downstream classifier
$G$ operating on a projection $\varphi$ is bounded above in its
safety-relevant information by $I(A;\varphi(X))$: the mutual information
between the admissibility variable and the representation. If the projection
has destroyed the admissibility information, no classifier can recover it.
The guardrail is interpellating a void. Its ``Hey, you there!'' reverberates
in an empty room.

Žižek's reading of the symptom becomes applicable here. The symptom, for Žižek,
is not a disguised expression of an underlying truth that interpretation would
reveal. It is a \emph{production} of the structure that generated the symptom
in order to sustain itself. The AI safety discourse---with its endless
elaboration of guardrails, red-teaming exercises, constitutional AI, RLHF,
output filters---is, from the admissibility perspective, a symptomatic
elaboration: the endless production of solutions to a problem that is
systematically displaced from its actual locus (the representation) to a
manageable but ultimately symptomatic site (the classifier). The symptom
persists not because the solutions are inadequate but because the problem
they address is not the real problem. And the real problem cannot be addressed
from within the discourse that the symptom sustains.

This is not a cynical observation. It has a precise technical correlate: the
Bayes-risk floor theorem establishes that when $I(A;\varphi(X)) \approx 0$,
the minimal achievable classification error approaches 50\%---the rate of
chance. The safety apparatus, operating on such a projection, performs at
the level of a coin flip. This is not a failure of implementation; it is the
symptom revealing, in the mode of the absurd, the foreclosure that the
discourse of safety has systematically refused.

% ══════════════════════════════════════════════════════════════════════════
\section{The Divided Subject of Planning:\\
Towards a Psychoanalysis of the Act}

The Lacanian subject is constitutively divided: barred ($\$$), split between the
symbolic order it inhabits and the Real that the symbolic order cannot absorb.
This division is not a deficiency to be overcome but the condition of possibility
of desire---and, more relevantly here, of the act. The act, for Lacan and
subsequently for Žižek, is the moment when the subject traverses the fantasy,
ceases to be supported by the \emph{objet a} that organises its desire, and
confronts the void directly.

The planning subject---whether artificial or biological---is constitutively
barred in exactly the admissibility-theoretic sense. It is barred by the
observational--interventional gap: it observes from within the symbolic order
of its representation, but it acts from a position that the representation
cannot fully capture. The reachability structure that guides action exceeds
the observational representation that learning produces. The planning subject
is always already $\$$: split between the representation it has and the
reachability structure it needs.

The \emph{act of planning}, in this frame, is the moment of traversal. The
agent acts, and in acting, produces evidence about the reachability structure
that its representation did not preserve. The trajectory-sampling diagnostic---
take two states the representation identifies, try actions from each,
observe whether outcomes diverge---is the formal analogue of the act as Lacan
describes it: a controlled intervention into the Real that forces the Real to
reveal itself through its effects. Witnessed admissibility distortion is the
symptom revealed by the act; the violation detected is the objet a surfacing
momentarily before the fantasy reconstitutes itself.

And here is the crucial point that the admissibility framework implicitly
but inexorably makes: the act is epistemically prior to the representation.
What the agent can do is not derived from what the agent represents; rather,
what the agent represents is evaluated by its adequacy to what the agent can
do. The admissibility relation is derived from the requirements of action
prior to any specification of reward, observation, or architecture. In
Lacanian terms: the Real of reachability is prior to the symbolic order of
the representation, and the symbolic order is evaluated by its (necessarily
incomplete) articulation of the Real.

This is the structural reason why the Bayesian parallel the paper develops
is illuminating. Bayesian inference, as \citet{jaynes2003probability} makes
explicit through the Knuth--Skilling derivation \citep{knuth2012foundations},
is not an arbitrary methodology but the unique consistent extension of logic
to degrees of plausibility. Similarly, the admissibility relation is not one
criterion among many but the unique coarsest safe quotient for planning. Both
are derived from structural consistency requirements rather than positive
empirical content. Both name a Real that the symbolic order must approximate
without ever coinciding with.

% ══════════════════════════════════════════════════════════════════════════
\section{Coda: Repetition, Difference, and the\\
Infinite Deferral of the Admissible}

Deleuze \citeyearpar{deleuze1994difference} distinguishes between the repetition
of the Same---the empty repetition that produces no difference---and the
repetition of Difference itself: the return that transforms what it repeats.
The research programme that the admissibility framework calls for---develop
estimators for $D_A(\varphi)$, design admissibility-preserving objectives,
audit representations for witnessed violations---is not a repetition of the Same.
It does not return to the predictive paradigm and attempt to do it better.
It returns to the fundamental question of representation---what must be preserved?---
and repeats it \emph{as Difference}: on the other side of the
observational--interventional gap, from the enunciative position that the
predictive paradigm forecloses.

The admissibility quotient $\mathcal{X}/{\sim_A}$ is, in Deleuzian terms, an
\emph{Idea} in the technical sense: a differential structure that organises a
problem space without resolving into a determinate solution. It specifies what
must be preserved without specifying how preservation is to be achieved. It is
immanent to the problem of planning---derived from what planning requires---
without being transcendent to any particular implementation. It is the virtual
structure that any actual representation must actualise, inadequately, partially,
always with remainder.

The remainder is $D_A(\varphi)$. And the remainder, as every reader of Lacan
knows, never disappears. It circulates. It insists. It returns in the mode of
the symptom---as planning failure, as safety classifier collapse, as the
endless proliferation of guardrails that cannot guard against what the
representation has already foreclosed. The symptom is the message. And the
message is: \emph{there is no rapport d'observation que tienne.}

The suture will not hold.

\bigskip
\noindent\textbf{Acknowledgements.} The author has no institutional affiliation
and no conflicts of interest to declare.

% ══════════════════════════════════════════════════════════════════════════
\begin{thebibliography}{99}

\bibitem[Althusser(1971)]{althusser1971ideology}
Althusser, L. (1971).
Ideology and ideological state apparatuses.
In \textit{Lenin and Philosophy and Other Essays}.
Monthly Review Press.

\bibitem[Badiou(1988)]{badiou1988being}
Badiou, A. (1988).
\textit{L'être et l'événement}.
Seuil. [\textit{Being and Event}, trans.\ O.\ Feltham, 2005, Continuum.]

\bibitem[Blum \& Blum(2022)]{blum2022conscious}
Blum, L., \& Blum, M. (2022).
A theoretical computer science perspective on consciousness.
\textit{Journal of Artificial Intelligence and Consciousness}, 9(1), 1--42.

\bibitem[Craik(1943)]{craik1943nature}
Craik, K.~J.~W. (1943).
\textit{The Nature of Explanation}.
Cambridge University Press.

\bibitem[Deleuze(1994)]{deleuze1994difference}
Deleuze, G. (1994).
\textit{Difference and Repetition}, trans.\ P.\ Patton.
Columbia University Press. [Original: \textit{Différence et répétition}, 1968.]

\bibitem[Deleuze \& Guattari(1987)]{deleuze1987thousand}
Deleuze, G., \& Guattari, F. (1987).
\textit{A Thousand Plateaus: Capitalism and Schizophrenia}, trans.\ B.\ Massumi.
University of Minnesota Press.
[Original: \textit{Mille plateaux}, 1980.]

\bibitem[Gibson(1979)]{gibson1979ecological}
Gibson, J.~J. (1979).
\textit{The Ecological Approach to Visual Perception}.
Houghton Mifflin.

\bibitem[Jaynes(2003)]{jaynes2003probability}
Jaynes, E.~T. (2003).
\textit{Probability Theory: The Logic of Science}.
Cambridge University Press.

\bibitem[Knuth \& Skilling(2012)]{knuth2012foundations}
Knuth, K.~H., \& Skilling, J. (2012).
Foundations of inference.
\textit{Axioms}, 1(1), 38--73.

\bibitem[Lacan(1973)]{lacan1973four}
Lacan, J. (1973).
\textit{Le Séminaire, Livre XI: Les quatre concepts fondamentaux de la
psychanalyse}.
Seuil. [\textit{The Four Fundamental Concepts of Psychoanalysis},
trans.\ A.\ Sheridan, 1978, Norton.]

\bibitem[Lacan(1975)]{lacan1975encore}
Lacan, J. (1975).
\textit{Le Séminaire, Livre XX: Encore}.
Seuil. [\textit{On Feminine Sexuality, the Limits of Love and Knowledge},
trans.\ B.\ Fink, 1998, Norton.]

\bibitem[Pearl(2009)]{pearl2009causality}
Pearl, J. (2009).
\textit{Causality: Models, Reasoning, and Inference}, 2nd ed.
Cambridge University Press.

\bibitem[Poincar\'e(1914)]{poincare1914}
Poincar\'e, H. (1914).
\textit{Science and Method}.
Thomas Nelson.

\bibitem[Varela et al.(1999)]{varela1999first}
Varela, F.~J., Shear, J., et al.\ (Eds.) (1999).
\textit{The View from Within: First-Person Approaches to the Study
of Consciousness}.
Imprint Academic.

\bibitem[Žižek(1989)]{zizek1989sublime}
Žižek, S. (1989).
\textit{The Sublime Object of Ideology}.
Verso.

\bibitem[Žižek(1999)]{zizek1999ticklish}
Žižek, S. (1999).
\textit{The Ticklish Subject: The Absent Centre of Political Ontology}.
Verso.

\bibitem[Žižek(2006)]{zizek2006parallax}
Žižek, S. (2006).
\textit{The Parallax View}.
MIT Press.

\end{thebibliography}

\end{document}
